\documentclass[11pt,a4paper]{article}
\usepackage[utf8]{inputenc}
\usepackage[T1]{fontenc}
\usepackage[catalan]{babel}
\usepackage{lmodern}
\usepackage[sfdefault,lining]{FiraSans}
\usepackage[a4paper,top=2.3cm,bottom=2.2cm,left=1.8cm,right=1.8cm,headheight=1.3cm,headsep=0.4cm,footskip=1cm]{geometry}
\usepackage{amsmath,amssymb}
\usepackage{icomma}
\usepackage{graphicx}
\usepackage[table]{xcolor}
\usepackage{tikz}
\usepackage{fontawesome5}
\usepackage{ccicons}
\usepackage{fancyhdr}
\usepackage{lastpage}
\usepackage{tasks}
\usepackage{needspace}
\usepackage[most]{tcolorbox}
\usepackage{array}
\usepackage[hidelinks,pdftitle={Fitxa d'exercicis de la unitat 9: Perímetres i àrees. Teorema de Pitàgores. 3r d'ESO},pdfsubject={Matemàtiques 3r ESO},pdfkeywords={teorema de Pitàgores, hipotenusa, catets, perímetre, àrea, polígons regulars, apotema, circumferència, cercle, sector circular, corona circular}]{hyperref}

\definecolor{navy}{HTML}{1B2748}
\definecolor{accent}{HTML}{E8603A}
\definecolor{teal}{HTML}{1F9E90}
\definecolor{soft}{HTML}{EEF1F8}
\definecolor{muted}{HTML}{5D6685}

\setlength{\parindent}{0pt}
\setlength{\parskip}{3pt}
\settasks{label=\alph*), label-format=\bfseries\color{navy}, label-width=1.3em, item-indent=1.8em, column-sep=0.6em, after-item-skip=5pt, before-skip=3pt, after-skip=4pt}

% ---------- capçalera i peu ----------
\newcommand{\llicencia}{\href{https://creativecommons.org/licenses/by-nc-sa/4.0/deed.ca}{\ccbyncsa}\ \ Llicència Creative Commons Reconeixement-NoComercial-CompartirIgual 4.0 Internacional (CC BY-NC-SA 4.0)}
\pagestyle{fancy}
\fancyhf{}
\fancyhead[L]{\raisebox{-0.35\height}{\includegraphics[height=0.95cm]{logo-gva.png}}}
\fancyhead[C]{\small\color{muted}Fitxa d'exercicis · Unitat 9: Perímetres i àrees. Teorema de Pitàgores · 3r d'ESO}
\fancyhead[R]{\raisebox{-0.35\height}{\includegraphics[height=0.95cm]{logo-vermellar.png}}}
\fancyfoot[L]{\scriptsize\color{muted}\llicencia}
\fancyfoot[R]{\small\color{muted}\thepage}
\renewcommand{\headrulewidth}{0.4pt}
\renewcommand{\headrule}{\color{muted!40}\hrule height \headrulewidth}
\fancypagestyle{primera}{\fancyhf{}\renewcommand{\headrulewidth}{0pt}
  \fancyfoot[L]{\scriptsize\color{muted}\llicencia}
  \fancyfoot[R]{\small\color{muted}\thepage}}

% ---------- nivells i exercicis ----------
\newcommand{\estrelles}[1]{{\color{accent}\ifcase#1\or\faStar\faStar[regular]\faStar[regular]\or\faStar\faStar\faStar[regular]\or\faStar\faStar\faStar\fi}}
\newcounter{exer}
\newcommand{\ex}[1]{\par\needspace{4\baselineskip}\medskip\stepcounter{exer}\noindent{\color{navy}\bfseries\theexer.}\ \ {\small\estrelles{#1}}\ \ }

\newcommand{\seccio}[2]{\par\needspace{10\baselineskip}\bigskip
  \begin{tcolorbox}[colback=navy,colframe=navy,coltext=white,arc=2mm,boxrule=0pt,left=3mm,right=3mm,top=1.2mm,bottom=1.2mm]
  {\large\bfseries #1\quad #2}
  \end{tcolorbox}\vspace{-2pt}}

\newcolumntype{C}[1]{>{\centering\arraybackslash}p{#1}}
\usepackage{pifont}

% quadre de recordatori (només a la fitxa d'exercicis)
\newcommand{\recorda}[1]{\par\smallskip\begin{tcolorbox}[colback=accent!7,colframe=accent!55,arc=2mm,boxrule=0.5pt,left=2.5mm,right=2.5mm,top=1mm,bottom=1mm]\small{\color{accent}\faLightbulb[regular]}\ \textbf{Recorda:}\ #1\end{tcolorbox}}
\newcommand{\res}[1]{#1}
\newcommand{\rest}[1]{#1}

\def\fmneed{14}
% =====================================================================
% Capçalera FentMates de les fitxes: un color per a cada curs.
% Els colors són els de la pàgina inicial (LEVELS a assets/js/materials.js);
% el logo ve de _eines/logo/pdf/ (python3 _eines/logo/genera.py).
%
% Ús, al final del preàmbul de la unitat:
%   % =====================================================================
% Capçalera FentMates de les fitxes: un color per a cada curs.
% Els colors són els de la pàgina inicial (LEVELS a assets/js/materials.js);
% el logo ve de _eines/logo/pdf/ (python3 _eines/logo/genera.py).
%
% Ús, al final del preàmbul de la unitat:
%   % =====================================================================
% Capçalera FentMates de les fitxes: un color per a cada curs.
% Els colors són els de la pàgina inicial (LEVELS a assets/js/materials.js);
% el logo ve de _eines/logo/pdf/ (python3 _eines/logo/genera.py).
%
% Ús, al final del preàmbul de la unitat:
%   \input{../comu/fentmates.tex}
%   \fmcurs{1eso}
% i, al començament del document:
%   \capcalera{Matemàtiques}{1r ESO}{Fitxa d'exercicis · Unitat 2}{Divisibilitat}
% Opcional abans de \input: \def\fmneed{14} (espai mínim abans de cada secció).
% =====================================================================
\makeatletter
\def\fm@logodir{../../logo/pdf}
\providecommand{\fmneed}{10}
\def\fm@cursnom{}

% clau del curs, color i color del text damunt del color
\newcommand{\fm@curs}[3]{\@namedef{fm@color@#1}{#2}\@namedef{fm@text@#1}{#3}}
\fm@curs{1eso}{E6B84A}{navy}
\fm@curs{2eso}{3B82F6}{white}
\fm@curs{3eso}{FF7A50}{white}
\fm@curs{4eso-a}{A78BFA}{white}
\fm@curs{4eso-b}{6D28D9}{white}
\fm@curs{1batx-mi}{9B1C1C}{white}
\fm@curs{1batx-ccss}{EF6B6B}{white}
\fm@curs{2batx-mii}{1E6B47}{white}
\fm@curs{2batx-ccss}{2BB3A3}{white}

\newcommand{\fmcurs}[1]{%
  \@ifundefined{fm@color@#1}{\PackageError{fentmates}{Curs desconegut: #1}{}}{}%
  \edef\fm@tmp{\noexpand\definecolor{curs}{HTML}{\@nameuse{fm@color@#1}}%
               \noexpand\colorlet{curstext}{\@nameuse{fm@text@#1}}}\fm@tmp
  \colorlet{cursfosc}{curs!72!black}%
  \def\fm@logo{\fm@logodir/vertical-#1.pdf}%
  \def\fm@marca{\fm@logodir/marca-#1.pdf}}

% ---------- capçalera de la primera pàgina ----------
% #1 matèria (petit), #2 curs, #3 tipus de document i unitat, #4 títol
\newcommand{\capcalera}[4]{%
  \gdef\fm@cursnom{#2}%
  \par\noindent
  \begin{tcolorbox}[colback=curs,colframe=curs,boxrule=0pt,arc=2.2mm,width=3.5cm,height=2.15cm,
      valign=center,left=3mm,right=2mm,top=1mm,bottom=1mm,nobeforeafter,box align=center]
    \color{curstext}{\footnotesize #1}\par\vspace{1pt}{\fontsize{19}{21}\selectfont\bfseries #2\par}%
  \end{tcolorbox}\hspace{5mm}%
  \begin{minipage}[c]{\dimexpr\linewidth-3.5cm-5mm-2.9cm\relax}
    {\color{cursfosc}\small\bfseries\MakeUppercase{#3}}\par\vspace{2pt}
    {\color{navy}\fontsize{20}{23}\selectfont\bfseries #4\par}
  \end{minipage}\hfill
  \begin{minipage}[c]{2.6cm}\raggedleft\includegraphics[width=2.6cm]{\fm@logo}\end{minipage}\par
  \vspace{2pt}}

% ---------- barres de secció del color del curs ----------
\renewcommand{\seccio}[2]{\par\needspace{\fmneed\baselineskip}\bigskip
  \begin{tcolorbox}[colback=curs,colframe=curs,coltext=curstext,arc=2mm,boxrule=0pt,left=3mm,right=3mm,top=1.2mm,bottom=1.2mm]
  {\large\bfseries #1\quad #2}
  \end{tcolorbox}\vspace{-2pt}}

% ---------- capçalera i peu de les pàgines ----------
% primera pàgina: logos institucionals al peu
\fancypagestyle{primera}{\fancyhf{}\renewcommand{\headrulewidth}{0pt}
  \fancyfoot[L]{\raisebox{-0.32\height}{\includegraphics[height=0.78cm]{logo-gva.png}}\hspace{5mm}%
    \raisebox{-0.32\height}{\includegraphics[height=0.78cm]{logo-vermellar.png}}}
  \fancyfoot[R]{\scriptsize\color{muted}\href{https://creativecommons.org/licenses/by-nc-sa/4.0/deed.ca}{\ccbyncsa}\ \ CC BY-NC-SA 4.0 · fentmates.org\qquad\small\thepage}}
% la resta: el símbol a l'esquerra i el curs a la dreta, amb el títol al centre com abans
\fancyhead[L]{\raisebox{-0.3\height}{\includegraphics[height=0.62cm]{\fm@marca}}}
\fancyhead[R]{\small\bfseries\color{cursfosc}\fm@cursnom}
\renewcommand{\headrule}{\color{curs!45}\hrule height \headrulewidth}
\makeatother

%   \fmcurs{1eso}
% i, al començament del document:
%   \capcalera{Matemàtiques}{1r ESO}{Fitxa d'exercicis · Unitat 2}{Divisibilitat}
% Opcional abans de \input: \def\fmneed{14} (espai mínim abans de cada secció).
% =====================================================================
\makeatletter
\def\fm@logodir{../../logo/pdf}
\providecommand{\fmneed}{10}
\def\fm@cursnom{}

% clau del curs, color i color del text damunt del color
\newcommand{\fm@curs}[3]{\@namedef{fm@color@#1}{#2}\@namedef{fm@text@#1}{#3}}
\fm@curs{1eso}{E6B84A}{navy}
\fm@curs{2eso}{3B82F6}{white}
\fm@curs{3eso}{FF7A50}{white}
\fm@curs{4eso-a}{A78BFA}{white}
\fm@curs{4eso-b}{6D28D9}{white}
\fm@curs{1batx-mi}{9B1C1C}{white}
\fm@curs{1batx-ccss}{EF6B6B}{white}
\fm@curs{2batx-mii}{1E6B47}{white}
\fm@curs{2batx-ccss}{2BB3A3}{white}

\newcommand{\fmcurs}[1]{%
  \@ifundefined{fm@color@#1}{\PackageError{fentmates}{Curs desconegut: #1}{}}{}%
  \edef\fm@tmp{\noexpand\definecolor{curs}{HTML}{\@nameuse{fm@color@#1}}%
               \noexpand\colorlet{curstext}{\@nameuse{fm@text@#1}}}\fm@tmp
  \colorlet{cursfosc}{curs!72!black}%
  \def\fm@logo{\fm@logodir/vertical-#1.pdf}%
  \def\fm@marca{\fm@logodir/marca-#1.pdf}}

% ---------- capçalera de la primera pàgina ----------
% #1 matèria (petit), #2 curs, #3 tipus de document i unitat, #4 títol
\newcommand{\capcalera}[4]{%
  \gdef\fm@cursnom{#2}%
  \par\noindent
  \begin{tcolorbox}[colback=curs,colframe=curs,boxrule=0pt,arc=2.2mm,width=3.5cm,height=2.15cm,
      valign=center,left=3mm,right=2mm,top=1mm,bottom=1mm,nobeforeafter,box align=center]
    \color{curstext}{\footnotesize #1}\par\vspace{1pt}{\fontsize{19}{21}\selectfont\bfseries #2\par}%
  \end{tcolorbox}\hspace{5mm}%
  \begin{minipage}[c]{\dimexpr\linewidth-3.5cm-5mm-2.9cm\relax}
    {\color{cursfosc}\small\bfseries\MakeUppercase{#3}}\par\vspace{2pt}
    {\color{navy}\fontsize{20}{23}\selectfont\bfseries #4\par}
  \end{minipage}\hfill
  \begin{minipage}[c]{2.6cm}\raggedleft\includegraphics[width=2.6cm]{\fm@logo}\end{minipage}\par
  \vspace{2pt}}

% ---------- barres de secció del color del curs ----------
\renewcommand{\seccio}[2]{\par\needspace{\fmneed\baselineskip}\bigskip
  \begin{tcolorbox}[colback=curs,colframe=curs,coltext=curstext,arc=2mm,boxrule=0pt,left=3mm,right=3mm,top=1.2mm,bottom=1.2mm]
  {\large\bfseries #1\quad #2}
  \end{tcolorbox}\vspace{-2pt}}

% ---------- capçalera i peu de les pàgines ----------
% primera pàgina: logos institucionals al peu
\fancypagestyle{primera}{\fancyhf{}\renewcommand{\headrulewidth}{0pt}
  \fancyfoot[L]{\raisebox{-0.32\height}{\includegraphics[height=0.78cm]{logo-gva.png}}\hspace{5mm}%
    \raisebox{-0.32\height}{\includegraphics[height=0.78cm]{logo-vermellar.png}}}
  \fancyfoot[R]{\scriptsize\color{muted}\href{https://creativecommons.org/licenses/by-nc-sa/4.0/deed.ca}{\ccbyncsa}\ \ CC BY-NC-SA 4.0 · fentmates.org\qquad\small\thepage}}
% la resta: el símbol a l'esquerra i el curs a la dreta, amb el títol al centre com abans
\fancyhead[L]{\raisebox{-0.3\height}{\includegraphics[height=0.62cm]{\fm@marca}}}
\fancyhead[R]{\small\bfseries\color{cursfosc}\fm@cursnom}
\renewcommand{\headrule}{\color{curs!45}\hrule height \headrulewidth}
\makeatother

%   \fmcurs{1eso}
% i, al començament del document:
%   \capcalera{Matemàtiques}{1r ESO}{Fitxa d'exercicis · Unitat 2}{Divisibilitat}
% Opcional abans de \input: \def\fmneed{14} (espai mínim abans de cada secció).
% =====================================================================
\makeatletter
\def\fm@logodir{../../logo/pdf}
\providecommand{\fmneed}{10}
\def\fm@cursnom{}

% clau del curs, color i color del text damunt del color
\newcommand{\fm@curs}[3]{\@namedef{fm@color@#1}{#2}\@namedef{fm@text@#1}{#3}}
\fm@curs{1eso}{E6B84A}{navy}
\fm@curs{2eso}{3B82F6}{white}
\fm@curs{3eso}{FF7A50}{white}
\fm@curs{4eso-a}{A78BFA}{white}
\fm@curs{4eso-b}{6D28D9}{white}
\fm@curs{1batx-mi}{9B1C1C}{white}
\fm@curs{1batx-ccss}{EF6B6B}{white}
\fm@curs{2batx-mii}{1E6B47}{white}
\fm@curs{2batx-ccss}{2BB3A3}{white}

\newcommand{\fmcurs}[1]{%
  \@ifundefined{fm@color@#1}{\PackageError{fentmates}{Curs desconegut: #1}{}}{}%
  \edef\fm@tmp{\noexpand\definecolor{curs}{HTML}{\@nameuse{fm@color@#1}}%
               \noexpand\colorlet{curstext}{\@nameuse{fm@text@#1}}}\fm@tmp
  \colorlet{cursfosc}{curs!72!black}%
  \def\fm@logo{\fm@logodir/vertical-#1.pdf}%
  \def\fm@marca{\fm@logodir/marca-#1.pdf}}

% ---------- capçalera de la primera pàgina ----------
% #1 matèria (petit), #2 curs, #3 tipus de document i unitat, #4 títol
\newcommand{\capcalera}[4]{%
  \gdef\fm@cursnom{#2}%
  \par\noindent
  \begin{tcolorbox}[colback=curs,colframe=curs,boxrule=0pt,arc=2.2mm,width=3.5cm,height=2.15cm,
      valign=center,left=3mm,right=2mm,top=1mm,bottom=1mm,nobeforeafter,box align=center]
    \color{curstext}{\footnotesize #1}\par\vspace{1pt}{\fontsize{19}{21}\selectfont\bfseries #2\par}%
  \end{tcolorbox}\hspace{5mm}%
  \begin{minipage}[c]{\dimexpr\linewidth-3.5cm-5mm-2.9cm\relax}
    {\color{cursfosc}\small\bfseries\MakeUppercase{#3}}\par\vspace{2pt}
    {\color{navy}\fontsize{20}{23}\selectfont\bfseries #4\par}
  \end{minipage}\hfill
  \begin{minipage}[c]{2.6cm}\raggedleft\includegraphics[width=2.6cm]{\fm@logo}\end{minipage}\par
  \vspace{2pt}}

% ---------- barres de secció del color del curs ----------
\renewcommand{\seccio}[2]{\par\needspace{\fmneed\baselineskip}\bigskip
  \begin{tcolorbox}[colback=curs,colframe=curs,coltext=curstext,arc=2mm,boxrule=0pt,left=3mm,right=3mm,top=1.2mm,bottom=1.2mm]
  {\large\bfseries #1\quad #2}
  \end{tcolorbox}\vspace{-2pt}}

% ---------- capçalera i peu de les pàgines ----------
% primera pàgina: logos institucionals al peu
\fancypagestyle{primera}{\fancyhf{}\renewcommand{\headrulewidth}{0pt}
  \fancyfoot[L]{\raisebox{-0.32\height}{\includegraphics[height=0.78cm]{logo-gva.png}}\hspace{5mm}%
    \raisebox{-0.32\height}{\includegraphics[height=0.78cm]{logo-vermellar.png}}}
  \fancyfoot[R]{\scriptsize\color{muted}\href{https://creativecommons.org/licenses/by-nc-sa/4.0/deed.ca}{\ccbyncsa}\ \ CC BY-NC-SA 4.0 · fentmates.org\qquad\small\thepage}}
% la resta: el símbol a l'esquerra i el curs a la dreta, amb el títol al centre com abans
\fancyhead[L]{\raisebox{-0.3\height}{\includegraphics[height=0.62cm]{\fm@marca}}}
\fancyhead[R]{\small\bfseries\color{cursfosc}\fm@cursnom}
\renewcommand{\headrule}{\color{curs!45}\hrule height \headrulewidth}
\makeatother

\fmcurs{3eso}

\begin{document}
\thispagestyle{primera}

\capcalera{Matemàtiques}{3r ESO}{Fitxa d'exercicis · Unitat 9}{Perímetres i àrees. Teorema de Pitàgores}
\vspace{10pt}
Nom i cognoms: \rule{8.4cm}{0.4pt}\hfill Grup: \rule{1.6cm}{0.4pt}\hfill Data: \rule{2.2cm}{0.4pt}
\vspace{8pt}

\begin{tcolorbox}[colback=soft,colframe=soft,arc=2mm,left=3mm,right=3mm,top=1.5mm,bottom=1.5mm]
\small Exercicis graduats per nivells de dificultat:\quad \estrelles{1}\ bàsic \qquad \estrelles{2}\ mitjà \qquad \estrelles{3}\ avançat.\\[2pt] Resol-los al quadern mostrant tots els passos i fes un dibuix de cada situació. Si un resultat no és exacte, arrodoneix-lo a les centèsimes; per a $\pi$, fes servir la tecla de la calculadora.
\end{tcolorbox}

% =====================================================
\seccio{1}{El teorema de Pitàgores}

\ex{1} Calcula la hipotenusa d'un triangle rectangle els catets del qual mesuren (si no és exacta, arrodoneix a les centèsimes):\begin{tasks}(3)
  \task 6 cm i 8 cm
  \task 5 cm i 12 cm
  \task 9 m i 40 m
  \task 3 cm i 7 cm
  \task 1,5 m i 2 m
  \task 10 cm i 10 cm
\end{tasks}\recorda{En un triangle rectangle, $c^2=a^2+b^2$, on $c$ és la \textbf{hipotenusa} (el costat oposat a l'angle recte, el més llarg) i $a$, $b$ són els \textbf{catets}. Per a la hipotenusa sumem: $c=\sqrt{a^2+b^2}$; per a un catet restem: $a=\sqrt{c^2-b^2}$.}

\ex{1} Calcula el catet que falta en cada triangle rectangle ($c$ és la hipotenusa):\begin{tasks}(3)
  \task $c=25$, $b=7$
  \task $c=17$, $a=8$
  \task $c=10$, $b=8$
  \task $c=12$, $b=5$
  \task $c=6{,}5$, $a=2{,}5$
  \task $c=9$, $a=4$
\end{tasks}

\ex{1} Calcula la longitud $x$ en cada triangle rectangle (mesures en cm):\par\smallskip\begin{center}\begin{tikzpicture}[x=0.15cm,y=0.15cm,baseline=(current bounding box.center),line join=round]
\draw[fill=accent!22,draw=accent,thick] (0.000,0.000) -- (12.000,0.000) -- (0.000,16.000) -- cycle;
\draw[thin] (1.467,0.000) -- (1.467,1.467) -- (0.000,1.467);
\node[font=\small,below] at (6.000,0.000) {12};
\node[font=\small,left] at (0.000,8.000) {16};
\node[font=\small,above right] at (6.000,8.000) {$x$};
\end{tikzpicture}\qquad \begin{tikzpicture}[x=0.15cm,y=0.15cm,baseline=(current bounding box.center),line join=round]
\draw[fill=accent!22,draw=accent,thick] (0.000,0.000) -- (12.000,0.000) -- (0.000,5.000) -- cycle;
\draw[thin] (1.467,0.000) -- (1.467,1.467) -- (0.000,1.467);
\node[font=\small,below] at (6.000,0.000) {$x$};
\node[font=\small,left] at (0.000,2.500) {5};
\node[font=\small,above right] at (6.000,2.500) {13};
\end{tikzpicture}\qquad \begin{tikzpicture}[x=0.15cm,y=0.15cm,baseline=(current bounding box.center),line join=round]
\draw[fill=accent!22,draw=accent,thick] (0.000,0.000) -- (15.000,0.000) -- (0.000,8.000) -- cycle;
\draw[thin] (1.467,0.000) -- (1.467,1.467) -- (0.000,1.467);
\node[font=\small,below] at (7.500,0.000) {15};
\node[font=\small,left] at (0.000,4.000) {$x$};
\node[font=\small,above right] at (7.500,4.000) {17};
\end{tikzpicture}\end{center}

\ex{2} Una escala de 6,5 m està recolzada en una paret. El peu de l'escala està a 2,5 m de la paret. A quina altura de la paret arriba?

\ex{2} Volem un estel amb un fil de 50 m tens. Si estem a 30 m (en horitzontal) del punt de terra que hi ha just davall de l'estel, a quina altura vola? (No tingues en compte l'altura de la mà.)

\ex{2} Un vaixell ix del port i navega 9 km cap al nord i, després, 12 km cap a l'est. A quina distància en línia recta es troba del port?

\ex{2} Una rampa salva un desnivell de 0,8 m i ocupa 6 m en horitzontal. Quant mesura la rampa?

\ex{3} En un triangle rectangle, un catet mesura 2 cm més que l'altre i la hipotenusa mesura 10 cm. Quant mesuren els catets?

\ex{3} Un pal vertical de 9 m es trenca pel vent i la punta toca a terra a 3 m de la base (la part de baix continua dreta). A quina altura s'ha trencat?


% =====================================================
\seccio{2}{Classificació de triangles pels costats}

\ex{1} Classifica segons els angles (acutangle, rectangle o obtusangle) els triangles de costats:\begin{tasks}(3)
  \task 5, 12 i 13
  \task 6, 7 i 8
  \task 4, 6 i 9
  \task 10, 24 i 26
  \task 7, 8 i 12
  \task 9, 10 i 11
\end{tasks}\recorda{Si $c$ és el costat més llarg: $c^2<a^2+b^2$ \ $\Rightarrow$ acutangle; \ $c^2=a^2+b^2$ \ $\Rightarrow$ rectangle; \ $c^2>a^2+b^2$ \ $\Rightarrow$ obtusangle.}

\ex{1} Quines d'aquestes ternes són pitagòriques?\begin{tasks}(5)
  \task 12, 16, 20
  \task 9, 12, 16
  \task 15, 36, 39
  \task 11, 60, 61
  \task 8, 9, 12
\end{tasks}

\ex{2} Poden ser els costats d'un triangle? Si és així, classifica'l.\begin{tasks}(3)
  \task 3, 5 i 9
  \task 4, 4 i 7
  \task 5, 7 i 12
\end{tasks}

\ex{2} Un triangle té els costats d'1,5 m, 2 m i 2,5 m. És rectangle? Si ho és, calcula'n l'àrea prenent els catets com a base i altura.

\ex{2} Un fuster vol comprovar si una prestatgeria de 90 cm × 120 cm té els angles rectes. Quant ha de mesurar la diagonal? Si la mesura i li dona 152 cm, com és l'angle?

\ex{3} Quant ha de mesurar $x$ perquè el triangle de costats 9, 12 i $x$ siga rectangle? Hi ha dues possibilitats.

\ex{3} El triangle té costats 6, 8 i $x$, amb $x$ un nombre enter més gran que 8. Per a quins valors de $x$ és acutangle? I rectangle? I obtusangle?


% =====================================================
\seccio{3}{Aplicacions del teorema de Pitàgores}

\ex{1} Calcula la diagonal de:\begin{tasks}(3)
  \task un rectangle de 20 × 15 cm
  \task un quadrat de 7 cm de costat
  \task una porta de 2,4 × 1 m
\end{tasks}\recorda{Busca el triangle rectangle amagat: la diagonal d'un rectangle, l'altura d'un triangle isòsceles o equilàter, la meitat de les diagonals d'un rombe o l'altura d'un trapezi en formen un. Decideix quin costat és la hipotenusa i aplica Pitàgores.}

\ex{1} Calcula l'altura d'un triangle isòsceles de costats iguals de 17 cm i base de 16 cm, i la d'un triangle equilàter de 12 cm de costat.

\ex{1} Les diagonals d'un rombe mesuren 18 cm i 24 cm. Calcula'n el costat i el perímetre.

\ex{2} La diagonal d'un quadrat mesura 12 cm. Calcula'n el costat i el perímetre.

\ex{2} Un trapezi rectangle té les bases de 15 cm i 9 cm i l'altura de 8 cm. Calcula el costat oblic i el perímetre.

\ex{2} Un trapezi isòsceles té les bases de 22 cm i 12 cm i els costats no paral·lels de 13 cm. Calcula'n l'altura.

\ex{2} Calcula la distància entre els punts:\begin{tasks}(3)
  \task $A(2,\,1)$ i $B(5,\,5)$
  \task $C(-3,\,4)$ i $D(3,\,-4)$
  \task $E(0,\,-2)$ i $F(4,\,3)$
\end{tasks}\recorda{$d(A,\,B)=\sqrt{(x_2-x_1)^2+(y_2-y_1)^2}$: és la hipotenusa d'un triangle rectangle de catets la diferència de les $x$ i la de les $y$.}

\ex{3} Siga el triangle de vèrtexs $A(1,\,1)$, $B(7,\,1)$ i $C(4,\,5)$. Calcula'n els costats i el perímetre, i digues com és segons els costats i segons els angles.

\ex{3} Un rombe té 10 cm de costat i una diagonal de 12 cm. Quant mesura l'altra diagonal?

\ex{3} Un televisor de 32 polzades té la pantalla de proporció 16:9 (l'amplària i l'altura són $16k$ i $9k$). Les polzades mesuren la diagonal i 1 polzada són 2,54 cm. Calcula l'amplària i l'altura de la pantalla.


% =====================================================
\seccio{4}{Perímetres i àrees de triangles i quadrilàters}

\ex{1} Calcula l'àrea i el perímetre de:\begin{tasks}(2)
  \task un quadrat de 9 cm de costat
  \task un rectangle de 12 cm × 5 cm
  \task un romboide de base 10 cm, altura 6 cm i costat oblic 7 cm
  \task un rombe de diagonals 16 cm i 12 cm
\end{tasks}\recorda{Quadrat $A=l^2$; \ rectangle i romboide $A=b\cdot h$; \ rombe $A=\frac{D\cdot d}{2}$; \ trapezi $A=\frac{(B+b)\cdot h}{2}$; \ triangle $A=\frac{b\cdot h}{2}$. El perímetre és la suma dels costats.}

\ex{1} Calcula l'àrea de:\begin{tasks}(2)
  \task un trapezi de bases 11 m i 7 m i altura 5 m
  \task un triangle de 14 cm de base i 9 cm d'altura
\end{tasks}

\ex{1} Expressa en la unitat que s'indica:\begin{tasks}(4)
  \task $3{,}5\text{ m}^2$ en cm²
  \task $45\,000\text{ cm}^2$ en m²
  \task $2{,}3\text{ ha}$ en m²
  \task $750\text{ dm}^2$ en m²
\end{tasks}

\ex{1} Un rectangle té 72 cm² d'àrea i 9 cm de base. Calcula'n l'altura i el perímetre.

\ex{2} Calcula l'àrea d'un triangle equilàter de 8 cm de costat.

\ex{2} Calcula l'àrea i el perímetre d'un triangle isòsceles de costats iguals de 10 cm i base de 12 cm.

\ex{2} Calcula l'àrea i el perímetre d'un trapezi isòsceles de bases 20 cm i 10 cm i costats no paral·lels de 13 cm.

\ex{2} Un rombe té 40 cm de perímetre i una diagonal de 16 cm. Calcula'n l'àrea.

\ex{2} Un camp rectangular fa 120 m × 85 m. Expressa'n la superfície en m² i en hectàrees. Quant costarà sembrar-hi gespa a 3 € el m²?

\ex{3} Volem pintar una paret de 4,5 m × 2,6 m que té una finestra d'1,2 m × 1 m. Si un pot de pintura cobreix 4 m², quants pots hem de comprar?

\ex{3} Els catets d'un triangle rectangle mesuren 20 cm i 15 cm. Calcula'n l'àrea, la hipotenusa i l'altura sobre la hipotenusa.

\ex{3} Un rectangle té 34 cm de perímetre i 13 cm de diagonal. Quant mesuren els costats? Quina és l'àrea?


% =====================================================
\seccio{5}{Polígons regulars}

\ex{1} Calcula el perímetre i l'àrea dels polígons regulars:\begin{tasks}(1)
  \task pentàgon de 6 cm de costat i 4,13 cm d'apotema
  \task hexàgon de 5 cm de costat i 4,33 cm d'apotema
  \task octàgon de 4 cm de costat i 4,83 cm d'apotema
\end{tasks}\recorda{En un polígon regular, $A=\frac{P\cdot a}{2}$, on $P$ és el perímetre i $a$ l'apotema. El radi, l'apotema i mig costat formen un triangle rectangle: $r^2=a^2+\left(\frac l2\right)^2$. En l'hexàgon regular, $r=l$.}

\ex{1} Un hexàgon regular té 8 cm de costat. Quant mesura el radi? Calcula'n l'apotema i l'àrea.

\ex{2} Un pentàgon regular té 12 cm de costat i 10,2 cm de radi. Calcula'n l'apotema i l'àrea.

\ex{2} Un quadrat està inscrit en una circumferència de 6 cm de radi. Calcula'n el costat, l'apotema i l'àrea.

\ex{2} Volem cobrir 3 m² de paret amb rajoles en forma d'hexàgon regular de 10 cm de costat. Quantes rajoles calen, com a mínim?

\ex{3} Un triangle equilàter està inscrit en una circumferència de 4 cm de radi. Sabent que la seua apotema és la meitat del radi, calcula'n el costat i l'àrea.


% =====================================================
\seccio{6}{La circumferència: longitud i arcs}

\ex{1} Calcula la longitud de la circumferència de:\begin{tasks}(3)
  \task 7 cm de radi
  \task 15 cm de diàmetre
  \task 2,5 m de radi
\end{tasks}\recorda{$L=2\pi r=\pi d$. \ Un arc de $n^\circ$ mesura $L_{\text{arc}}=\frac{2\pi r\cdot n}{360}$. Fes servir la tecla $\pi$ de la calculadora i arrodoneix només al final.}

\ex{1} Una circumferència fa 62,83 cm de longitud. Quant mesuren el radi i el diàmetre?

\ex{1} Calcula la longitud dels arcs:\begin{tasks}(3)
  \task de 90° i 8 cm de radi
  \task de 30° i 12 cm de radi
  \task de 225° i 4 cm de radi
\end{tasks}

\ex{2} Les rodes d'una bicicleta fan 60 cm de diàmetre. Quantes voltes fa cada roda en un recorregut d'1,5 km?

\ex{2} La corda interior d'una pista d'atletisme està formada per dues rectes de 100 m i dues semicircumferències de 36,5 m de radi. Quant mesura una volta?

\ex{2} L'agulla minutera d'un rellotge fa 15 cm. Quina distància recorre la seua punta en 20 minuts?

\ex{3} En una circumferència de 12 cm de radi, un arc mesura 18,85 cm. Quants graus té?

\ex{3} La Terra fa una volta al Sol cada any seguint una òrbita gairebé circular de 150 milions de km de radi. A quina velocitat mitjana (en km/h) viatja la Terra?


% =====================================================
\seccio{7}{Àrea del cercle i figures circulars}

\ex{1} Calcula l'àrea del cercle de:\begin{tasks}(3)
  \task 5 cm de radi
  \task 12 cm de diàmetre
  \task 0,8 m de radi
\end{tasks}\recorda{Cercle $A=\pi r^2$; \ sector de $n^\circ$: $A=\frac{\pi r^2\cdot n}{360}$; \ corona: $A=\pi\left(R^2-r^2\right)$; \ trapezi circular: $A=\frac{\pi\left(R^2-r^2\right)\cdot n}{360}$.}

\ex{1} Calcula l'àrea dels sectors circulars:\begin{tasks}(3)
  \task de 60° i 6 cm de radi
  \task de 150° i 10 cm de radi
  \task de 270° i 2 cm de radi
\end{tasks}

\ex{1} Calcula l'àrea d'una corona circular de radis 9 cm i 4 cm.

\ex{2} Un cercle té 200 cm² d'àrea. Calcula'n el radi i la longitud de la circumferència.

\ex{2} Calcula l'àrea d'un trapezi circular de radis 12 cm i 7 cm i amplitud 80°.

\ex{2} Una pizza de 30 cm de diàmetre es talla en 8 porcions iguals. Calcula la superfície de cada porció i la longitud de la seua vora (l'arc).

\ex{2} Un hexàgon regular té 30 cm de perímetre. Calcula l'àrea del cercle circumscrit (el que passa pels vèrtexs) i la del cercle inscrit (el radi és l'apotema).

\ex{3} Calcula l'àrea del segment circular ombrejat, determinat per un angle de 60° en un cercle de 6 cm de radi. (El triangle és equilàter.)\par\smallskip\begin{center}\begin{tikzpicture}[x=0.55cm,y=0.55cm,baseline=(current bounding box.center),line join=round]
\draw[fill=navy!18,draw=navy,thick] (0.000,0.000) -- (3.000,0.000) arc[start angle=0,end angle=60,radius=3] -- cycle;
\draw[fill=white,draw=white] (0.000,0.000) -- (3.000,0.000) -- (1.500,2.598) -- cycle;
\draw[thick] (0.000,0.000) -- (3.000,0.000);
\draw[thick] (0.000,0.000) -- (1.500,2.598);
\draw[thick] (3.000,0.000) -- (1.500,2.598);
\draw[thick] (3.000,0.000) arc[start angle=0,end angle=60,radius=3];
\fill (0.000,0.000) circle (1.4pt);
\node[font=\small,below] at (1.500,0.000) {6 cm};
\draw[thick] (0.600,0.000) arc[start angle=0,end angle=60,radius=0.6];
\node[font=\small,right] at (0.750,0.350) {\scriptsize 60°};
\end{tikzpicture}\end{center}

\ex{3} Una corona circular té una àrea de $50\pi$ cm² i el radi exterior fa 10 cm. Quant mesura el radi interior?


% =====================================================
\seccio{8}{Figures compostes i problemes}

\ex{1} Calcula l'àrea i el perímetre de la façana de la caseta.\par\smallskip\begin{center}\begin{tikzpicture}[x=0.17cm,y=0.17cm,baseline=(current bounding box.center),line join=round]
\draw[fill=teal!12,draw=teal,thick] (0.000,0.000) -- (6.000,0.000) -- (6.000,6.000) -- (3.000,10.000) -- (0.000,6.000) -- cycle;
\draw[dashed,muted] (0.000,6.000) -- (6.000,6.000);
\draw[dashed,muted] (3.000,6.000) -- (3.000,10.000);
\draw[thin] (4.294,6.000) -- (4.294,7.294) -- (3.000,7.294);
\node[font=\small,below] at (3.000,0.000) {6 m};
\node[font=\small,right] at (6.000,3.000) {6 m};
\node[font=\small,right] at (3.000,8.000) {4 m};
\end{tikzpicture}\end{center}\recorda{Per a una figura composta, descompon-la en figures conegudes i suma les àrees, o veu-la com una figura gran a la qual en lleves trossos i resta. El perímetre és només el contorn exterior.}

\ex{1} A un quadrat de 8 cm de costat li afegim un semicercle sobre un dels costats. Calcula l'àrea i el perímetre de la figura.\par\smallskip\begin{center}\begin{tikzpicture}[x=0.15cm,y=0.15cm,baseline=(current bounding box.center),line join=round]
\draw[fill=navy!18,draw=navy,thick] (0.000,0.000) -- (8.000,0.000) -- (8.000,8.000) -- (0.000,8.000) -- cycle;
\draw[fill=navy!18,draw=navy,thick] (4.000,8.000) -- (8.000,8.000) arc[start angle=0,end angle=180,radius=4] -- cycle;
\draw[dashed,muted] (0.000,8.000) -- (8.000,8.000);
\node[font=\small,below] at (4.000,0.000) {8 cm};
\end{tikzpicture}\end{center}

\ex{2} D'un rectangle de 10 cm × 6 cm retallem dos semicercles de 6 cm de diàmetre en els costats curts. Calcula l'àrea i el perímetre de la figura que queda.\par\smallskip\begin{center}\begin{tikzpicture}[x=0.2cm,y=0.2cm,baseline=(current bounding box.center),line join=round]
\draw[fill=navy!18,draw=navy,thick] (0.000,0.000) -- (10.000,0.000) -- (10.000,6.000) -- (0.000,6.000) -- cycle;
\draw[fill=white,draw=white] (0.000,3.000) -- (0.000,0.000) arc[start angle=-90,end angle=90,radius=3] -- cycle;
\draw[fill=white,draw=white] (10.000,3.000) -- (10.000,6.000) arc[start angle=90,end angle=270,radius=3] -- cycle;
\draw[thick] (0.000,0.000) arc[start angle=-90,end angle=90,radius=3];
\draw[thick] (10.000,6.000) arc[start angle=90,end angle=270,radius=3];
\node[font=\small,below] at (5.000,0.000) {10 cm};
\node[font=\small,right] at (10.000,3.000) {6 cm};
\end{tikzpicture}\end{center}

\ex{2} Calcula l'àrea de la zona ombrejada: un quadrat de 12 cm de costat menys un quart de cercle de radi 12 cm.\par\smallskip\begin{center}\begin{tikzpicture}[x=0.13cm,y=0.13cm,baseline=(current bounding box.center),line join=round]
\draw[fill=navy!18,draw=navy,thick] (0.000,0.000) -- (12.000,0.000) -- (12.000,12.000) -- (0.000,12.000) -- cycle;
\draw[fill=white,draw=white] (0.000,0.000) -- (12.000,0.000) arc[start angle=0,end angle=90,radius=12] -- cycle;
\draw[thick] (12.000,0.000) arc[start angle=0,end angle=90,radius=12];
\node[font=\small,below] at (6.000,0.000) {12 cm};
\draw[dashed,muted] (0.000,0.000) -- (8.485,8.485);
\node[font=\small,above left] at (4.200,4.200) {12};
\end{tikzpicture}\end{center}

\ex{2} Una finestra té forma de rectangle d'1,4 m d'amplària i 1,8 m d'altura acabat en un semicercle. Calcula la superfície de vidre i la longitud del marc exterior.\par\smallskip\begin{center}\begin{tikzpicture}[x=0.75cm,y=0.75cm,baseline=(current bounding box.center),line join=round]
\draw[fill=teal!12,draw=teal,thick] (0.000,0.000) -- (1.400,0.000) -- (1.400,1.800) -- (0.000,1.800) -- cycle;
\draw[fill=teal!12,draw=teal,thick] (0.700,1.800) -- (1.400,1.800) arc[start angle=0,end angle=180,radius=0.7] -- cycle;
\draw[dashed,muted] (0.000,1.800) -- (1.400,1.800);
\node[font=\small,below] at (0.700,0.000) {1,4 m};
\node[font=\small,right] at (1.400,0.900) {1,8 m};
\end{tikzpicture}\end{center}

\ex{2} Al voltant d'una piscina de 25 m × 12 m es fa una vorera de 2 m d'amplària. Calcula la superfície de la vorera.

\ex{3} Calcula l'àrea de la «fulla» ombrejada. Els arcs són quarts de circumferència amb centre en dos vèrtexs oposats del quadrat de 10 cm de costat.\par\smallskip\begin{center}\begin{tikzpicture}[x=0.15cm,y=0.15cm,baseline=(current bounding box.center),line join=round]
\draw[fill=teal!12,draw=teal,thick] (0.000,0.000) -- (10.000,0.000) -- (10.000,10.000) -- (0.000,10.000) -- cycle;
\draw[fill=navy!18,draw=navy,thick] (10,0) arc[start angle=0,end angle=90,radius=10] arc[start angle=180,end angle=270,radius=10] -- cycle;
\node[font=\small,below] at (5.000,0.000) {10 cm};
\end{tikzpicture}\end{center}

\ex{3} Tres cercles iguals caben just dins d'un rectangle de 30 cm × 10 cm. Calcula l'àrea de la zona ombrejada.\par\smallskip\begin{center}\begin{tikzpicture}[x=0.14cm,y=0.14cm,baseline=(current bounding box.center),line join=round]
\draw[fill=navy!18,draw=navy,thick] (0.000,0.000) -- (30.000,0.000) -- (30.000,10.000) -- (0.000,10.000) -- cycle;
\draw[fill=white,draw=white] (5.000,5.000) circle[radius=5];
\draw[fill=white,draw=white] (15.000,5.000) circle[radius=5];
\draw[fill=white,draw=white] (25.000,5.000) circle[radius=5];
\draw[draw=teal,thick] (5.000,5.000) circle[radius=5];
\draw[draw=teal,thick] (15.000,5.000) circle[radius=5];
\draw[draw=teal,thick] (25.000,5.000) circle[radius=5];
\draw[dashed,muted] (0.000,0.000) -- (30.000,0.000) -- (30.000,10.000) -- (0.000,10.000) -- cycle;
\node[font=\small,below] at (15.000,0.000) {30 cm};
\node[font=\small,right] at (30.000,5.000) {10 cm};
\end{tikzpicture}\end{center}

\ex{3} Calcula l'àrea compresa entre un hexàgon regular de 4 cm de costat i la circumferència que passa pels seus vèrtexs.

\ex{3} Una cabra està lligada amb una corda de 5 m a una cantonada exterior d'un corral quadrat de 4 m de costat, i no pot entrar al corral. Quina superfície pot pasturar?\par\smallskip\begin{center}\begin{tikzpicture}[x=0.24cm,y=0.24cm,baseline=(current bounding box.center),line join=round]
\draw[fill=navy!18,draw=navy,thick] (0.000,0.000) -- (0.000,5.000) arc[start angle=90,end angle=360,radius=5] -- cycle;
\draw[fill=navy!18,draw=navy,thick] (4.000,0.000) -- (5.000,0.000) arc[start angle=0,end angle=90,radius=1] -- cycle;
\draw[fill=navy!18,draw=navy,thick] (0.000,4.000) -- (1.000,4.000) arc[start angle=0,end angle=90,radius=1] -- cycle;
\draw[fill=teal!12,draw=teal,thick] (0.000,0.000) -- (4.000,0.000) -- (4.000,4.000) -- (0.000,4.000) -- cycle;
\fill (0.000,0.000) circle (1.4pt);
\draw[dashed,muted] (0.000,0.000) -- (-3.540,-3.540);
\node[font=\small,above left] at (-1.800,-1.800) {5 m};
\node[font=\small,inner sep=1pt] at (2.000,2.000) {corral};
\node[font=\small,below] at (2.000,4.000) {4 m};
\end{tikzpicture}\end{center}

\vfill
\begin{center}\color{muted}\small\faCheckCircle[regular]\ \ Ànim! Si et costa algun exercici, repassa els apunts interactius de la unitat.\end{center}

\end{document}
